% GENERATED FILE. Edit canonical lessons, then run npm run generate:books.
% canonical-source-hash: fnv1a64:e4fdab8bea0477e5
% canonical-lessons: AR-C83-bati, AR-C83-sari, AR-C83-sahl, AR-C83-sab, AR-C83-rakhis

\chapter{Slow, Fast, Easy, Difficult, Cheap}
\label{ch:ar-slow-fast-easy-difficult-cheap}
\hlchaptermodality{83}

\begin{chapteropening}
\textbf{By the end of this chapter:} \emph{I can say slow, fast, easy, difficult and cheap.}

Complete the last lesson of chapter 83: I can say slow, fast, easy, difficult and cheap.
\end{chapteropening}

\section[\texorpdfstring{\ar{بطيء} (baṭīʾ)}{بطيء (baṭīʾ)}]{\ar{بطيء} (baṭīʾ) — slow}

\label{lesson:AR-C83-bati}



\begin{quote}
Before the new one: say the Arabic for long, then the Arabic for short.
\end{quote}

\subsection*{You'll want to know: \ar{بطيء}}
\textbf{\ar{بطيء}} — \emph{baṭīʾ} — \textquotedblleft{}slow\textquotedblright{}.

\begin{cousinweb}
\textbf{\ar{بطيء}} (\emph{baṭīʾ}) is built on the native Arabic root \emph{b-ṭ-ʾ}, \textquotedblleft{}to be slow\textquotedblright{} (Lane). \emph{Baṭīʾ} is the \emph{faʿīl} adjective, \textquotedblleft{}slow\textquotedblright{}.
\end{cousinweb}

\begin{grammarlens}[title={What you've built}]
One more word for this chapter.
\end{grammarlens}

\subsection*{Your turn}
\begin{itemize}
\raggedright
  \item \emph{Say it:} \emph{baṭīʾ}
  \item \emph{Say it:} \emph{baṭīʾ}, once more
  \item \emph{Say it:} say \emph{baṭīʾ}
  \item \emph{Recall:} say the Arabic for long, then the Arabic for short, then say \emph{baṭīʾ} again
\end{itemize}

\begin{tcolorbox}[breakable,colback=teal!4,colframe=teal!35!black,title={Before you move on}]
What does \ar{بطيء} mean? (\textquotedblleft{}Slow\textquotedblright{}.) Say it once more.
\end{tcolorbox}
\section[\texorpdfstring{\ar{سريع} (sarīʿ)}{سريع (sarīʿ)}]{\ar{سريع} (sarīʿ) — fast}

\label{lesson:AR-C83-sari}



\begin{quote}
Before the new one: say the Arabic for short, then the Arabic for slow.
\end{quote}

\subsection*{You'll want to know: \ar{سريع}}
\textbf{\ar{سريع}} — \emph{sarīʿ} — \textquotedblleft{}fast\textquotedblright{}.

\begin{cousinweb}
\textbf{\ar{سريع}} (\emph{sarīʿ}) is built on the native Arabic root \emph{s-r-ʿ}, \textquotedblleft{}to be quick\textquotedblright{} (Lane). \emph{Sarīʿ} is the \emph{faʿīl} adjective, \textquotedblleft{}fast\textquotedblright{}.
\end{cousinweb}

\begin{grammarlens}[title={What you've built}]
One more word for this chapter.
\end{grammarlens}

\subsection*{Your turn}
\begin{itemize}
\raggedright
  \item \emph{Say it:} \emph{sarīʿ}
  \item \emph{Say it:} \emph{sarīʿ}, once more
  \item \emph{Say it:} say \emph{sarīʿ}
  \item \emph{Recall:} say the Arabic for short, then the Arabic for slow, then say \emph{sarīʿ} again
\end{itemize}

\begin{tcolorbox}[breakable,colback=teal!4,colframe=teal!35!black,title={Before you move on}]
What does \ar{سريع} mean? (\textquotedblleft{}Fast\textquotedblright{}.) Say it once more.
\end{tcolorbox}
\section[\texorpdfstring{\ar{سهل} (sahl)}{سهل (sahl)}]{\ar{سهل} (sahl) — easy}

\label{lesson:AR-C83-sahl}



\begin{quote}
Before the new one: say the Arabic for slow, then the Arabic for fast.
\end{quote}

\subsection*{You'll want to know: \ar{سهل}}
\textbf{\ar{سهل}} — \emph{sahl} — \textquotedblleft{}easy\textquotedblright{}.

\begin{cousinweb}
\textbf{\ar{سهل}} (\emph{sahl}) is built on the native Arabic root \emph{s-h-l}, \textquotedblleft{}to be easy; level ground\textquotedblright{} (Lane). \emph{Sahl} is also level, plain ground.
\end{cousinweb}

\begin{grammarlens}[title={What you've built}]
One more word for this chapter.
\end{grammarlens}

\subsection*{Your turn}
\begin{itemize}
\raggedright
  \item \emph{Say it:} \emph{sahl}
  \item \emph{Say it:} \emph{sahl}, once more
  \item \emph{Say it:} say \emph{sahl}
  \item \emph{Recall:} say the Arabic for slow, then the Arabic for fast, then say \emph{sahl} again
\end{itemize}

\begin{tcolorbox}[breakable,colback=teal!4,colframe=teal!35!black,title={Before you move on}]
What does \ar{سهل} mean? (\textquotedblleft{}Easy\textquotedblright{}.) Say it once more.
\end{tcolorbox}
\section[\texorpdfstring{\ar{صعب} (ṣaʿb)}{صعب (ṣaʿb)}]{\ar{صعب} (ṣaʿb) — difficult}

\label{lesson:AR-C83-sab}



\begin{quote}
Before the new one: say the Arabic for fast, then the Arabic for easy.
\end{quote}

\subsection*{You'll want to know: \ar{صعب}}
\textbf{\ar{صعب}} — \emph{ṣaʿb} — \textquotedblleft{}difficult\textquotedblright{}.

\begin{cousinweb}
\textbf{\ar{صعب}} (\emph{ṣaʿb}) is built on the native Arabic root \emph{ṣ-ʿ-b}, \textquotedblleft{}to be hard, difficult\textquotedblright{} (Lane). \emph{Ṣaʿb} is the plain adjective, \textquotedblleft{}difficult\textquotedblright{}.
\end{cousinweb}

\begin{grammarlens}[title={What you've built}]
One more word for this chapter.
\end{grammarlens}

\subsection*{Your turn}
\begin{itemize}
\raggedright
  \item \emph{Say it:} \emph{ṣaʿb}
  \item \emph{Say it:} \emph{ṣaʿb}, once more
  \item \emph{Say it:} say \emph{ṣaʿb}
  \item \emph{Recall:} say the Arabic for fast, then the Arabic for easy, then say \emph{ṣaʿb} again
\end{itemize}

\begin{tcolorbox}[breakable,colback=teal!4,colframe=teal!35!black,title={Before you move on}]
What does \ar{صعب} mean? (\textquotedblleft{}Difficult\textquotedblright{}.) Say it once more.
\end{tcolorbox}
\section[\texorpdfstring{\ar{رخيص} (rakhīṣ)}{رخيص (rakhīṣ)}]{\ar{رخيص} (rakhīṣ) — cheap}

\label{lesson:AR-C83-rakhis}



\begin{quote}
Before the new one: say the Arabic for easy, then the Arabic for difficult.
\end{quote}

\subsection*{You'll want to know: \ar{رخيص}}
\textbf{\ar{رخيص}} — \emph{rakhīṣ} — \textquotedblleft{}cheap\textquotedblright{}.

\begin{cousinweb}
\textbf{\ar{رخيص}} (\emph{rakhīṣ}) is built on the native Arabic root \emph{r-kh-ṣ}, \textquotedblleft{}to be cheap; to be soft\textquotedblright{} (Lane). \emph{Rakhīṣ} is the \emph{faʿīl} adjective, \textquotedblleft{}cheap\textquotedblright{}.
\end{cousinweb}

\begin{grammarlens}[title={What you've built}]
One more word for this chapter.
\end{grammarlens}

\subsection*{Your turn}
\begin{itemize}
\raggedright
  \item \emph{Say it:} \emph{rakhīṣ}
  \item \emph{Say it:} \emph{rakhīṣ}, once more
  \item \emph{Say it:} say \emph{rakhīṣ}
  \item \emph{Recall:} say the Arabic for easy, then the Arabic for difficult, then say \emph{rakhīṣ} again
\end{itemize}

\begin{tcolorbox}[breakable,colback=teal!4,colframe=teal!35!black,title={Before you move on}]
What does \ar{رخيص} mean? (\textquotedblleft{}Cheap\textquotedblright{}.) Say it once more.
\end{tcolorbox}
